\documentclass[11pt]{article}
\usepackage[margin=1in]{geometry}
\usepackage{amsmath,amssymb,amsthm,hyperref}
\newtheorem{proposition}{Proposition}
\newcommand{\R}{\mathbb R}
\newcommand{\Al}{A_{\bullet}}
\newcommand{\Au}{A^{\bullet}}
\title{Transition convex hulls can differ for planar similarities}
\author{Principia Math in Progress}
\date{12 September 2026}
\begin{document}
\maketitle
\paragraph{Task 1. Resolve MathDB 354470.}
Conjecture 6.2 of Le\'sniak, Snigireva, Strobin and Vince
\cite{source}, p.~5425, proposes equality of the metric convex hulls
of the lower and upper transition attractors when each map takes
metrically convex sets onto metrically convex sets. We give a
counterexample on the Euclidean plane satisfying (H1)--(H3).

\begin{proposition}
Let $e=(1,0)$ and $J(x,y)=(-y,x)$. For $0\leq t\leq1$ set
\[
 r=r_t=\frac{1+t}{2},\qquad \varepsilon=1-r,
\]
and define the two-map IFS $F_t=\{f_t,g_t\}$ by
\begin{align*}
 f_t(x)&=rJ(x-\varepsilon e)+\varepsilon e,\\
 g_t(x)&=rJ^{-1}(x+\varepsilon e)-\varepsilon e.
\end{align*}
Its lower transition attractor is $\Al=\{0\}$. It has at least one
upper transition attractor, and every upper transition attractor
contains $(1,-1)$. Hence their metric convex hulls differ.
\end{proposition}

\begin{proof}
The maps are distinct invertible similarities for every parameter,
including both endpoints. Their coefficients depend continuously on
$t$, proving (H1), and their Lipschitz constants are $r<1$ for $t<1$,
proving (H2). Their respective unique fixed points are
$\varepsilon e$ and $-\varepsilon e$, both tending to zero, proving
(H3). Thus the limiting fixed-point set is $Q=\{0\}$.

In the Euclidean metric, metric segments are ordinary line segments.
Every invertible affine map $u$ satisfies
$u([x,y])=[u(x),u(y)]$, so both maps send every metrically convex set
onto a metrically convex set, as required in the conjecture.

At $t=1$, we have $F_1=\{J,J^{-1}\}$. The singleton $\{0\}$ is
invariant and contains $Q$, and every set containing $Q$ contains it.
The lower-transition definition therefore gives $\Al=\{0\}$.

We next verify that upper transition attractors exist. Since
$J^{-1}=-J$, the translations in the two affine maps are
\[
 b_+=\varepsilon(e-rJe),\qquad
 b_-=-\varepsilon(e+rJe).
\]
Their norms are $\varepsilon\sqrt{1+r^2}\leq2\varepsilon$.
Consequently the closed Euclidean ball $B=\{x:\|x\|\leq2\}$ is
invariant under both maps, because
\[
 \|f_t(x)\|,\ \|g_t(x)\|
 \leq2r+2(1-r)=2\qquad(x\in B).
\]
For $t<1$ the contraction theorem gives a nonempty compact attractor
$A_t\subset B$. The space of nonempty compact subsets of the compact
set $B$ is compact in the Hausdorff metric. Starting, for example,
with $t_n=1-1/n$, $n\geq1$, we may pass to a subsequence along which
$A_{t_n}$ converges. This gives an upper transition attractor as in
Definition 4.3 of the source.

To distinguish every such limit from the lower transition set,
consider the composition $P_t=f_t\circ g_t$. Direct calculation gives
\begin{equation}\label{composition}
 P_t(x)=r^2x+\varepsilon\bigl((1+r^2)e-2rJe\bigr).
\end{equation}
Its unique fixed point, for $t<1$, is
\begin{equation}\label{fixed}
 p_t=\frac{(1+r^2)e-2rJe}{1+r},
 \qquad \lim_{t\uparrow1}p_t=e-Je=(1,-1).
\end{equation}
For any $a\in A_t$, invariance gives $P_t^n(a)\in A_t$ for every
$n$. Since $P_t$ is a contraction, these points converge to $p_t$;
closedness gives $p_t\in A_t$.

Now let $\Au=\lim_n A_{t_n}$ be any upper transition attractor.
Hausdorff convergence and~\eqref{fixed} imply
\[
 \operatorname{dist}((1,-1),\Au)
 \leq\|(1,-1)-p_{t_n}\|+h(A_{t_n},\Au)\longrightarrow0.
\]
As $\Au$ is closed, $(1,-1)\in\Au$. In the source's notation,
$\operatorname{conv}_d S=\bigcup_{x,y\in S}[x,y]$. In particular
it contains $S$, because $[x,x]=\{x\}$. Hence
\[
 \operatorname{conv}_d\Al=\{0\},\qquad
 (1,-1)\in\operatorname{conv}_d\Au,
\]
which disproves the conjectured equality.
\end{proof}

The same formulas extend to all $t\geq0$. For $t>1$, $f_t$ expands
distances from its fixed point $\varepsilon e$ by the factor $r>1$,
so a generic orbit is unbounded. Thus $F_t$ is not contractive in any
equivalent complete metric, and cannot have an attractor with an open
basin. At $t=1$, $J^4$ is the identity, also precluding contraction
in any metric. The finite contractivity threshold required in
Section 4 is exactly one.

The centers vary with the parameter. Thus the construction meets the
general conjecture while lying outside the older semilinear
fixed-center hypothesis. It also works if one replaces the source's
union-of-segments hull by the ordinary convex hull or its closure.

\paragraph{Task 2. Determine the complete limiting set.}
Establish whether this family's attractors have a unique Hausdorff
limit, and calculate that limit explicitly.

\paragraph{Task 3. Quantify the effect of moving fixed points.}
For opposite rotations, characterize which rates and directions of
center motion produce nonzero limiting fixed points of compositions.

\paragraph{Task 4. Identify sufficient hypotheses for hull equality.}
Determine which restrictions on moving centers extend the known
fixed-center results to more general similarity families.

\paragraph{Task 5. Compare metric convexity conventions.}
Study equality of transition hulls defined by one union of segments,
iterated segment closure and closed metric convexity in non-Euclidean
spaces, under suitable additional hypotheses.

\begin{thebibliography}{1}
\bibitem{source} K. Le\'sniak, N. Snigireva, F. Strobin and A. Vince,
\emph{Transition phenomena for the attractor of an iterated function
system}, Nonlinearity \textbf{35} (2022), 5396--5426.
\href{https://doi.org/10.1088/1361-6544/ac8af1}{doi:10.1088/1361-6544/ac8af1}.
\end{thebibliography}
\end{document}
